\documentclass[a4paper,12pt]{report}

% ─── Packages ──────────────────────────────────────────────────────
\usepackage[utf8]{inputenc}
\usepackage[T1]{fontenc}
\usepackage[italian]{babel}
\usepackage{geometry}
\geometry{margin=2.5cm}
\usepackage{graphicx}
\usepackage{hyperref}
\hypersetup{
    colorlinks=true,
    linkcolor=blue!70!black,
    citecolor=blue!70!black,
    urlcolor=blue!60!black,
    pdfauthor={},
    pdftitle={Agentic UniBG – Relazione Finale},
}
\usepackage{enumitem}
\usepackage{tabularx}
\usepackage{booktabs}
\usepackage{longtable}
\usepackage{listings}
\usepackage{xcolor}
\usepackage{titlesec}
\usepackage{fancyhdr}
\usepackage{float}
\usepackage{caption}
\usepackage{microtype}
\usepackage{parskip}

% ─── Listings style ───────────────────────────────────────────────
\definecolor{codebg}{HTML}{F5F5F5}
\definecolor{codegreen}{HTML}{2E7D32}
\definecolor{codeblue}{HTML}{1565C0}
\definecolor{codegray}{HTML}{757575}
\definecolor{codepurple}{HTML}{7B1FA2}

\lstdefinestyle{codestyle}{
    backgroundcolor=\color{codebg},
    commentstyle=\color{codegreen}\itshape,
    keywordstyle=\color{codeblue}\bfseries,
    numberstyle=\tiny\color{codegray},
    stringstyle=\color{codepurple},
    basicstyle=\ttfamily\small,
    breakatwhitespace=false,
    breaklines=true,
    captionpos=b,
    keepspaces=true,
    numbers=left,
    numbersep=5pt,
    showspaces=false,
    showstringspaces=false,
    showtabs=false,
    tabsize=2,
    frame=single,
    rulecolor=\color{codegray!30},
}
\lstset{style=codestyle}

% ─── Headers / footers ────────────────────────────────────────────
\pagestyle{fancy}
\fancyhf{}
\fancyhead[L]{\small\leftmark}
\fancyhead[R]{\small Agentic UniBG}
\fancyfoot[C]{\thepage}
\renewcommand{\headrulewidth}{0.4pt}

% ─── Title formatting ─────────────────────────────────────────────
\titleformat{\chapter}[display]
  {\normalfont\huge\bfseries}{\chaptertitlename\ \thechapter}{18pt}{\Huge}
\titlespacing*{\chapter}{0pt}{-20pt}{30pt}

% ═══════════════════════════════════════════════════════════════════
\begin{document}

% ─── Title page ────────────────────────────────────────────────────
\begin{titlepage}
\centering
\vspace*{2cm}

{\Huge\bfseries Agentic UniBG\par}
\vspace{0.5cm}
{\Large Sistema Multi-Agent per l'Assistenza Universitaria\par}
\vspace{1.5cm}
{\large Relazione Finale del Progetto\par}
\vspace{0.5cm}
{\large Metodologia AMDD -- Agile Model Driven Development\par}
\vspace{3cm}

{\large Università degli Studi di Bergamo\par}
\vspace{0.5cm}
{\large Anno Accademico 2025/2026\par}

\vfill
{\small Documento generato il \today\par}
\end{titlepage}

% ─── Table of contents ─────────────────────────────────────────────
\tableofcontents
\newpage

% ═══════════════════════════════════════════════════════════════════
% CHAPTER 1 – INTRODUZIONE
% ═══════════════════════════════════════════════════════════════════
\chapter{Introduzione}

\section{Contesto del Progetto}

Il sito web dell'Università degli Studi di Bergamo (\texttt{unibg.it}) presenta una struttura gerarchica complessa, con informazioni distribuite su numerose pagine e documenti. Studenti e utenti devono navigare manualmente tra molte sezioni o consultare documenti lunghi per ricercare una singola informazione, con conseguenti perdite di tempo.

\textbf{Agentic UniBG} nasce con l'obiettivo di sviluppare un sistema intelligente multi-agent in grado di assistere studenti e utenti nella navigazione del sito universitario, fornendo risposte rapide, contestualizzate e personalizzate a qualsiasi domanda pertinente.

\section{Obiettivi del Sistema}

Il sistema è concepito come un assistente conversazionale avanzato, capace di:
\begin{itemize}
    \item Comprendere richieste in linguaggio naturale
    \item Classificare automaticamente l'intento della domanda
    \item Recuperare informazioni reali e aggiornate dal sito dell'ateneo
    \item Generare risposte strutturate, accurate e contestualizzate
    \item Personalizzare l'esperienza in base al profilo dello studente
    \item Gestire conversazioni multi-turno con memoria del contesto
\end{itemize}

\section{Metodologia di Sviluppo: AMDD}

Il progetto adotta la metodologia \textbf{Agile Model Driven Development (AMDD)}, che combina i principi dello sviluppo agile con la modellazione iterativa e incrementale.

L'approccio AMDD prevede:
\begin{enumerate}
    \item \textbf{Iterazione 0 -- Envisioning}: definizione della visione iniziale, dei requisiti ad alto livello e dell'architettura target. Nessun codice prodotto.
    \item \textbf{Iterazioni successive}: sviluppo incrementale con cicli brevi, ciascuno dei quali produce un incremento funzionante del sistema. I modelli UML vengono creati e aggiornati ad ogni iterazione secondo necessità (``just enough modeling'').
    \item \textbf{Evoluzione continua}: i requisiti, l'architettura e la documentazione evolvono di pari passo con il codice, senza fasi rigide di design upfront.
\end{enumerate}

In linea con i principi AMDD:
\begin{itemize}
    \item L'\textbf{Iterazione 0} ha prodotto la visione d'insieme, le user stories iniziali e l'architettura target
    \item Le \textbf{Iterazioni 1--4} hanno realizzato progressivamente le funzionalità, aggiornando i modelli (diagrammi UML, specifiche dei requisiti) ad ogni incremento
    \item L'\textbf{Iterazione 5} è attualmente \textit{work in progress}
\end{itemize}

\section{Tecnologie Utilizzate}

\begin{table}[H]
\centering
\begin{tabularx}{\textwidth}{l X}
\toprule
\textbf{Componente} & \textbf{Tecnologia} \\
\midrule
Frontend & React (Vite), servito tramite Nginx \\
Backend & Python, FastAPI, Uvicorn \\
Framework AI & LangChain, LangGraph \\
LLM Provider & Google Gemini 2.5 Flash (precedentemente Groq / LLaMA 3.3 70B) \\
Ricerca Web & Tavily API (search + extract) \\
Gestione PDF & PyMuPDF (fitz) \\
Database & MongoDB Atlas (Motor async driver) \\
Autenticazione & JWT (python-jose), bcrypt, httpOnly cookie \\
Containerizzazione & Docker, Docker Compose \\
\bottomrule
\end{tabularx}
\caption{Stack tecnologico del progetto}
\end{table}

\section{Struttura del Documento}

Il documento è organizzato seguendo la sequenza temporale delle iterazioni:
\begin{itemize}
    \item \textbf{Capitolo 2}: Iterazione 0 -- Visione iniziale e analisi preliminare
    \item \textbf{Capitolo 3}: Iterazione 1 -- MVP della pipeline multi-agent
    \item \textbf{Capitolo 4}: Iterazione 2 -- Autenticazione, profilazione e personalizzazione
    \item \textbf{Capitolo 5}: Iterazione 3 -- JWT, gestione sessione e memoria conversazionale
    \item \textbf{Capitolo 6}: Iterazione 4 -- Web search, date esami e pipeline logging
    \item \textbf{Capitolo 7}: Iterazione 5 -- Gestione profilo utente e miglioramenti (WIP)
    \item \textbf{Capitolo 8}: Scelte architetturali e di progetto
    \item \textbf{Capitolo 9}: Conclusioni e sviluppi futuri
\end{itemize}


% ═══════════════════════════════════════════════════════════════════
% CHAPTER 2 – ITERAZIONE 0
% ═══════════════════════════════════════════════════════════════════
\chapter{Iterazione 0 -- Visione Iniziale}

\section{Obiettivo dell'Iterazione}

L'Iterazione 0 rappresenta la fase di \textit{envisioning} prevista dalla metodologia AMDD. Non è stato prodotto codice funzionante, ma sono stati definiti:
\begin{itemize}
    \item La visione complessiva del progetto
    \item Le user stories iniziali (target completo)
    \item I requisiti funzionali e non funzionali ad alto livello
    \item L'architettura target del sistema
    \item I rischi iniziali identificati
\end{itemize}

\section{Visione del Progetto}

Agentic UniBG è stato concepito come un sistema multi-agent in grado di:
\begin{itemize}
    \item Supportare sia studenti autenticati che ospiti
    \item Comprendere domande in linguaggio naturale
    \item Accedere a fonti informative interne (database corsi, orari) ed esterne (web)
    \item Personalizzare l'esperienza in base al profilo studente
    \item Fornire risposte strutturate con citazione delle fonti
\end{itemize}

\section{User Stories Iniziali}

Le user stories della visione iniziale coprono l'intero spettro delle funzionalità desiderate:

\begin{table}[H]
\centering
\begin{tabularx}{\textwidth}{l X}
\toprule
\textbf{ID} & \textbf{Descrizione} \\
\midrule
US-01 & Come studente autenticato, voglio accedere al sistema con il mio profilo \\
US-02 & Come visitatore esterno, voglio porre domande senza autenticazione \\
US-03 & Come studente, voglio chiedere informazioni su un corso \\
US-04 & Come studente, voglio conoscere l'orario delle lezioni \\
US-05 & Come studente, voglio sapere se un'aula è libera \\
US-06 & Come studente, voglio trovare un regolamento o documento ufficiale \\
US-07 & Come studente, voglio essere guidato passo-passo in una procedura \\
US-08 & Come studente autenticato, voglio che il sistema ricordi le conversazioni \\
US-09 & Come utente, voglio fornire un feedback sulla risposta \\
\bottomrule
\end{tabularx}
\caption{User stories -- visione completa (Iterazione 0)}
\end{table}

\section{Architettura Target}

L'architettura iniziale prevede tre macro-componenti:

\begin{enumerate}
    \item \textbf{Frontend}: applicazione React per l'interazione utente (login, chat)
    \item \textbf{Backend (LangChain)}: API Gateway, Orchestrator Agent e agenti specializzati (Classifier, Profile, Memory, Web, Documents, Generator, Revision, Feedback)
    \item \textbf{Persistence Layer}: database SQL per i profili studente, database NoSQL per la cronologia delle conversazioni
    \item \textbf{Componenti esterni}: LLM Provider (inizialmente previsto Ollama), accesso Web/Internet
\end{enumerate}

Il flusso teorico completo prevede: l'utente invia una richiesta al Frontend, che la inoltra all'API Gateway; l'Orchestrator Agent la distribuisce agli agenti specializzati nell'ordine opportuno; la risposta finale risale all'Orchestrator e viene restituita all'utente.

\section{Use Case -- Visione Completa}

Sono stati identificati 13 macro-casi d'uso, suddivisi per attore:

\begin{table}[H]
\centering
\small
\begin{tabularx}{\textwidth}{l l X}
\toprule
\textbf{ID} & \textbf{Attore} & \textbf{Nome} \\
\midrule
UC-01 & Studente & Login come Studente \\
UC-02 & Studente & Registrazione Studente \\
UC-03 & Studente & Gestione Profilo \\
UC-04 & Ospite & Accesso come Ospite \\
UC-05 & Studente, Ospite & Inviare Domanda/Richiesta \\
UC-06 & Studente & Richiedere Guida Procedurale \\
UC-07 & Studente & Recupero Documento Ufficiale \\
UC-08 & Studente & Informazioni su un Corso \\
UC-09 & Studente & Consultazione Orari Lezioni \\
UC-10 & Studente & Verifica Occupazione Aule \\
UC-11 & Studente, Ospite & Informazioni Generali \\
UC-12 & Studente, Ospite & Ricevere Risposta \\
UC-13 & Studente, Ospite & Fornire Feedback \\
\bottomrule
\end{tabularx}
\caption{Use case della visione completa (Iterazione 0)}
\end{table}

\section{Rischi Iniziali e Strategia}

\begin{itemize}
    \item \textbf{Complessità architetturale elevata}: mitigata dalla scomposizione in iterazioni incrementali
    \item \textbf{Dipendenza da servizi LLM esterni}: mitigata dalla modularità (sostituzione provider trasparente)
    \item \textbf{Incertezza sulla qualità delle risposte}: mitigata dall'introduzione del Revision Agent e delle fonti web
\end{itemize}

La strategia adottata prevede di non implementare l'intera visione iniziale in un'unica fase, ma di procedere per \textit{minimum viable product} incrementali, evolvendo l'architettura progressivamente.


% ═══════════════════════════════════════════════════════════════════
% CHAPTER 3 – ITERAZIONE 1
% ═══════════════════════════════════════════════════════════════════
\chapter{Iterazione 1 -- MVP della Pipeline Multi-Agent}

\section{Obiettivo dell'Iterazione}

L'Iterazione 1 realizza il primo \textbf{Minimum Viable Product (MVP)} funzionante, limitato alla pipeline conversazionale di base:
\begin{enumerate}
    \item \textbf{Input Management}: ricezione della domanda dall'utente
    \item \textbf{Intent Classification}: classificazione della richiesta tramite Classifier Agent
    \item \textbf{Response Generation}: generazione della risposta tramite Generator Agent e LLM
    \item \textbf{Review}: revisione della risposta tramite Revision Agent
    \item \textbf{Output}: restituzione del risultato finale all'utente
\end{enumerate}

\section{User Stories Incluse}

\begin{table}[H]
\centering
\begin{tabularx}{\textwidth}{l X}
\toprule
\textbf{ID} & \textbf{Descrizione} \\
\midrule
US-10 & Come studente, voglio inviare una domanda testuale e ricevere una risposta \\
US-11 & Come sistema, voglio classificare la richiesta per generare una risposta coerente \\
US-12 & Come sistema, voglio generare una risposta tramite LLM pertinente alla domanda \\
US-13 & Come sistema, voglio verificare la qualità della risposta per ridurre errori \\
\bottomrule
\end{tabularx}
\caption{User stories -- Iterazione 1}
\end{table}

\section{Architettura Implementata}

L'architettura è organizzata in \textbf{due tier}:
\begin{itemize}
    \item \textbf{Client Tier}: interfaccia utente minimale per l'invio di query
    \item \textbf{Application Logic Tier}: Web Server (FastAPI) e App Server che ospitano la logica degli agenti
\end{itemize}

Non è presente alcun Data Tier, né autenticazione, registrazione o profilazione utente.

I componenti sviluppati sono:
\begin{itemize}
    \item \textbf{OrchestratorAgent}: coordina l'esecuzione degli agenti, costruisce il grafo LangGraph con tre nodi (\texttt{classify} $\rightarrow$ \texttt{generate} $\rightarrow$ \texttt{revise}), espone \texttt{handleQuery(query)}
    \item \textbf{ClassifierAgent}: classifica la query in una categoria semantica (es.\ \texttt{informazioni\_corso}, \texttt{orari}, \texttt{procedure}, \texttt{servizi}, \texttt{generale}, \texttt{altro}), espone \texttt{classify(query)}
    \item \textbf{GeneratorAgent}: genera la risposta utilizzando un prompt specifico per categoria, espone \texttt{generate(query, classification)}
    \item \textbf{RevisionAgent}: revisiona la bozza per correttezza e chiarezza, espone \texttt{revise(draft)}
    \item \textbf{AgentState}: dizionario tipizzato (\texttt{TypedDict}) che mantiene lo stato condiviso della pipeline
\end{itemize}

\section{Pipeline di Elaborazione}

Il flusso sequenziale implementato è:

\begin{verbatim}
User Input → Orchestrator → ClassifierAgent → GeneratorAgent 
           → RevisionAgent → Orchestrator → User Output
\end{verbatim}

\begin{enumerate}
    \item L'utente invia la query all'OrchestratorAgent
    \item L'Orchestrator invia la query al ClassifierAgent per la classificazione
    \item Il ClassifierAgent restituisce la categoria (intent label)
    \item L'Orchestrator invoca il GeneratorAgent con query e categoria
    \item Il GeneratorAgent produce la bozza di risposta
    \item L'Orchestrator invia la bozza al RevisionAgent
    \item Il RevisionAgent restituisce la risposta revisionata
    \item L'Orchestrator restituisce la risposta finale all'utente
\end{enumerate}

\section{Algoritmo di Classificazione}

L'algoritmo di classificazione utilizza un prompt template strutturato:
\begin{enumerate}
    \item Ricezione dell'input utente
    \item Costruzione di un prompt con l'elenco delle categorie disponibili
    \item Invocazione dell'LLM (Groq / LLaMA 3.3 70B in questa iterazione)
    \item Parsing dell'output: se la categoria restituita non è valida, fallback a \texttt{altro}
\end{enumerate}

\section{Requisiti Funzionali}

\begin{itemize}
    \item \textbf{UC-I1 -- Invio Richiesta}: l'utente invia una domanda tramite \texttt{POST /api/agent/query}
    \item \textbf{UC-I2 -- Classificazione Intento}: il ClassifierAgent analizza la query e la assegna a una categoria
    \item \textbf{UC-I3 -- Generazione Risposta}: il GeneratorAgent produce la risposta con prompt contestualizzato
    \item \textbf{UC-I4 -- Revisione Risposta}: il RevisionAgent revisiona la bozza per correttezza e chiarezza
\end{itemize}

\section{Requisiti Non Funzionali}

\begin{itemize}
    \item \textbf{RNF-I1}: tempo di risposta $\leq$ 10 secondi
    \item \textbf{RNF-I2}: modularità -- ogni agente è indipendente e sostituibile
    \item \textbf{RNF-I3}: robustezza -- fallimento del RevisionAgent non blocca la pipeline (graceful fallback)
    \item \textbf{RNF-I4}: statelessness -- nessuna persistenza tra richieste
    \item \textbf{RNF-I5}: estensibilità -- aggiunta di nuovi nodi al workflow senza modifiche alla logica esistente
    \item \textbf{RNF-I6}: documentazione API tramite Swagger UI (\texttt{/docs})
\end{itemize}

\section{Limitazioni}

\begin{itemize}
    \item Nessun database e nessuna persistenza
    \item Nessuna autenticazione o gestione utenti
    \item Nessuna personalizzazione delle risposte
    \item Nessun accesso a dati reali dell'ateneo (risposte solo da LLM)
    \item Possibili allucinazioni del modello
\end{itemize}


% ═══════════════════════════════════════════════════════════════════
% CHAPTER 4 – ITERAZIONE 2
% ═══════════════════════════════════════════════════════════════════
\chapter{Iterazione 2 -- Autenticazione, Profilazione e Personalizzazione}

\section{Obiettivo dell'Iterazione}

L'Iterazione 2 evolve il sistema dal prototipo dell'Iterazione 1 verso un'applicazione completa e deployabile, introducendo:
\begin{itemize}
    \item \textbf{Autenticazione utente}: registrazione e login con matricola e password
    \item \textbf{Profilazione personalizzata}: corso, dipartimento, anno, tipologia
    \item \textbf{Personalizzazione delle risposte}: gli agenti adattano le risposte al profilo
    \item \textbf{Frontend React}: interfaccia di login e chat completa
    \item \textbf{Persistenza su MongoDB}: tramite Motor (driver asincrono)
    \item \textbf{Containerizzazione}: Docker e Docker Compose
\end{itemize}

\section{User Stories Incluse}

\begin{table}[H]
\centering
\begin{tabularx}{\textwidth}{l X l}
\toprule
\textbf{ID} & \textbf{Descrizione} & \textbf{Ruolo} \\
\midrule
US-01 & Registrazione con matricola e password & Studente \\
US-02 & Autenticazione e ricezione profilo & Studente \\
US-03 & Invio domande senza autenticazione & Ospite \\
US-04 & Risposte personalizzate in base a corso e anno & Studente \\
US-05 & Invio richiesta e ricezione risposta elaborata & Utente \\
\bottomrule
\end{tabularx}
\caption{User stories -- Iterazione 2}
\end{table}

\section{Architettura a Tre Tier}

L'architettura evolve a \textbf{tre tier}:

\begin{table}[H]
\centering
\begin{tabularx}{\textwidth}{l X}
\toprule
\textbf{Tier} & \textbf{Tecnologia} \\
\midrule
Client Tier & React App (Vite + Nginx) \\
Backend Tier & FastAPI + Agenti (LangChain / LangGraph) \\
Data Tier & MongoDB Atlas (Motor async) \\
\bottomrule
\end{tabularx}
\end{table}

\section{Autenticazione e Registrazione}

\subsection{Registrazione}
\begin{enumerate}
    \item Lo studente compila il form (nome, cognome, matricola, password, dipartimento, corso, tipologia, anno)
    \item Il frontend invia \texttt{POST /api/auth/register}
    \item Il backend verifica l'unicità della matricola su MongoDB
    \item La password viene hashata con \texttt{bcrypt.hashpw} (con salt)
    \item Il documento utente viene salvato su MongoDB
    \item Il profilo viene restituito al frontend
\end{enumerate}

\subsection{Login}
\begin{enumerate}
    \item Lo studente inserisce matricola e password
    \item Il frontend invia \texttt{POST /api/auth/login}
    \item Il backend cerca lo studente per matricola su MongoDB
    \item \texttt{bcrypt.checkpw} verifica la password contro l'hash memorizzato
    \item In caso di successo, il profilo completo viene restituito
    \item Il frontend salva il profilo in memoria e abilita la chat
\end{enumerate}

\textbf{Nota importante}: in questa iterazione il \textbf{JWT non è stato implementato}. Il login restituisce direttamente i dati del profilo senza token di sessione. La gestione del token stateless è stata pianificata e rimandataa all'Iterazione 3.

\section{Personalizzazione delle Risposte}

La funzione \texttt{build\_user\_context(state)} costruisce una stringa di contesto dal profilo utente, iniettata nei prompt di tutti gli agenti:

\begin{itemize}
    \item Se l'utente è \textbf{autenticato}: include nome, cognome, matricola, dipartimento, corso, tipologia, anno
    \item Se l'utente è \textbf{ospite}: inserisce un disclaimer per risposte generiche
\end{itemize}

\section{Persistenza -- MongoDB}

Il modulo \texttt{db.py} utilizza Motor per la connessione asincrona a MongoDB Atlas. Il documento utente contiene: \texttt{name}, \texttt{surname}, \texttt{matricola}, \texttt{passwordHash}, \texttt{department}, \texttt{course}, \texttt{tipology}, \texttt{year}.

\section{Frontend React}

\begin{itemize}
    \item \textbf{LoginPage.jsx}: form di login/registrazione con selezione di dipartimento e corso da liste predefinite
    \item \textbf{App.jsx}: interfaccia chat principale con invio query e visualizzazione risposte
\end{itemize}

\section{Containerizzazione}

\begin{table}[H]
\centering
\begin{tabularx}{\textwidth}{l l l}
\toprule
\textbf{Servizio} & \textbf{Immagine} & \textbf{Porta} \\
\midrule
\texttt{frontend} & node + nginx & 3000 (→ 5173) \\
\texttt{backend} & python:3.11-slim & 8000 \\
\bottomrule
\end{tabularx}
\caption{Servizi Docker Compose}
\end{table}

\section{Endpoint API}

\begin{table}[H]
\centering
\small
\begin{tabularx}{\textwidth}{l l X}
\toprule
\textbf{Metodo} & \textbf{Path} & \textbf{Descrizione} \\
\midrule
GET & \texttt{/api/health} & Health check \\
POST & \texttt{/api/auth/login} & Login con matricola e password \\
POST & \texttt{/api/auth/register} & Registrazione nuovo studente \\
POST & \texttt{/api/agent/query} & Elaborazione query con pipeline agenti \\
GET & \texttt{/api/agents} & Lista agenti disponibili \\
GET & \texttt{/api/conversation/history} & Cronologia conversazioni \\
DELETE & \texttt{/api/conversation/history} & Cancella cronologia \\
POST & \texttt{/api/agent/analyze} & Analisi query senza esecuzione (debug) \\
\bottomrule
\end{tabularx}
\caption{Endpoint API -- Iterazione 2}
\end{table}

\section{Limitazioni}

\begin{itemize}
    \item \textbf{JWT non implementato}: il profilo è mantenuto lato frontend in memoria, senza token di sessione
    \item Nessun accesso a dati reali dell'ateneo
    \item Conversation history in memoria (non persistita)
    \item CORS con \texttt{allow\_origins=["*"]}
\end{itemize}


% ═══════════════════════════════════════════════════════════════════
% CHAPTER 5 – ITERAZIONE 3
% ═══════════════════════════════════════════════════════════════════
\chapter{Iterazione 3 -- JWT, Gestione Sessione e Memoria Conversazionale}

\section{Obiettivo dell'Iterazione}

L'Iterazione 3 consolida il sistema con tre macro-funzionalità:
\begin{itemize}
    \item \textbf{JWT Authentication via httpOnly Cookie}: il token di sessione viene generato lato server e trasmesso tramite cookie httpOnly, eliminando la gestione manuale del token nel frontend
    \item \textbf{Gestione della Sessione Browser}: ripristino automatico della sessione al ricaricamento della pagina tramite verifica del cookie JWT
    \item \textbf{Memoria Conversazionale}: la cronologia viene mantenuta nel \texttt{localStorage} del browser e trasmessa ad ogni query come \texttt{conversation\_history}
\end{itemize}

\section{User Stories Incluse}

\begin{table}[H]
\centering
\begin{tabularx}{\textwidth}{l X}
\toprule
\textbf{ID} & \textbf{Descrizione} \\
\midrule
US-06 & Il sistema ricorda la sessione dopo il ricaricamento della pagina \\
US-07 & Le risposte tengono conto delle domande precedenti \\
US-08 & Le conversazioni precedenti vengono ripristinate alla riapertura \\
US-09 & Logout sicuro con cancellazione della sessione \\
\bottomrule
\end{tabularx}
\caption{User stories -- Iterazione 3}
\end{table}

\section{Confronto con l'Iterazione 2}

\begin{table}[H]
\centering
\small
\begin{tabularx}{\textwidth}{l X X}
\toprule
\textbf{Funzionalità} & \textbf{Iterazione 2} & \textbf{Iterazione 3} \\
\midrule
JWT & Assente & httpOnly cookie con HMAC-SHA256 \\
Ripristino sessione & Assente & \texttt{GET /api/auth/verify} \\
Logout & Assente & \texttt{POST /api/auth/logout} + delete cookie \\
Memoria conversazionale & Assente & localStorage + conversation\_history \\
Auth Layer & Logica in main.py & Pacchetto \texttt{auth/} con 4 classi separate \\
CORS & \texttt{allow\_origins=["*"]} & Origini ristrette a localhost \\
\bottomrule
\end{tabularx}
\caption{Evoluzione dall'Iterazione 2 alla 3}
\end{table}

\section{Pacchetto \texttt{auth/} -- Refactoring}

Il codice di autenticazione è stato estratto in un pacchetto dedicato con responsabilità ben separate:

\begin{itemize}
    \item \textbf{JWTManager}: generazione (\texttt{generateToken}) e validazione (\texttt{validateToken}, \texttt{validateFromRequest}) dei JWT. Firma HMAC-SHA256, scadenza configurabile (default 7 giorni)
    \item \textbf{AuthService}: logica applicativa (\texttt{authenticate}, \texttt{createUser}) con verifica credenziali, hashing password e generazione token
    \item \textbf{AuthController}: ponte tra route handler FastAPI e AuthService; gestisce l'impostazione (\texttt{\_set\_auth\_cookie}) e l'eliminazione del cookie; filtra i campi sensibili (\texttt{\_public\_profile})
    \item \textbf{ProfileRepository}: operazioni CRUD sulla collection MongoDB (\texttt{findById}, \texttt{save}, \texttt{exists})
\end{itemize}

\section{Gestione della Sessione}

\subsection{Ripristino Automatico}
\begin{enumerate}
    \item Al caricamento dell'app React, viene effettuata \texttt{GET /api/auth/verify} con \texttt{credentials: "include"}
    \item Il backend estrae il JWT dall'httpOnly cookie tramite \texttt{JWTManager.validateFromRequest}
    \item Decodifica il payload, recupera la matricola e interroga MongoDB
    \item Restituisce il profilo completo al frontend
    \item Il frontend carica anche la conversazione da \texttt{localStorage}
\end{enumerate}

\subsection{Logout Sicuro}
\begin{enumerate}
    \item L'utente clicca logout
    \item \texttt{POST /api/auth/logout} cancella il cookie \texttt{access\_token} lato server
    \item Il frontend rimuove profilo, conversazione e reindirizza al login
\end{enumerate}

\section{Memoria Conversazionale}

\subsection{Persistenza in localStorage}
La cronologia è salvata con chiave \texttt{conversation\_\{matricola\}}:
\begin{itemize}
    \item Salvataggio ad ogni aggiornamento (React \texttt{useEffect})
    \item Ripristino dopo il verify token con successo
    \item Cancellazione al logout
\end{itemize}

\subsection{Trasmissione agli Agenti}
Ad ogni query, il frontend costruisce un array \texttt{conversation\_history} con gli ultimi \textbf{10 messaggi} (5 turni domanda/risposta) e lo invia nel body di \texttt{POST /api/agent/query}.

L'\texttt{AgentState} include il campo \texttt{conversation\_history: Optional[List[Dict]]}. L'Orchestrator lo distribuisce a:
\begin{itemize}
    \item \textbf{GeneratorAgent}: inietta la history come messaggi LLM nativi (\texttt{HumanMessage}/\texttt{AIMessage}) per vera memoria conversazionale
    \item \textbf{RevisionAgent}: considera la history nella revisione per evitare contraddizioni
\end{itemize}

\section{Sicurezza}

\begin{table}[H]
\centering
\begin{tabularx}{\textwidth}{l X X}
\toprule
\textbf{Aspetto} & \textbf{Iterazione 2} & \textbf{Iterazione 3} \\
\midrule
Token JWT & Non implementato & httpOnly cookie (inaccessibile da JS) \\
CORS & \texttt{allow\_origins=["*"]} & Origini esplicite (localhost) \\
Password & bcrypt & bcrypt (invariato) \\
Scadenza token & -- & 7 giorni (configurabile) \\
SameSite & -- & \texttt{lax} (protezione CSRF base) \\
\bottomrule
\end{tabularx}
\caption{Evoluzione sicurezza}
\end{table}

\section{Alternative di Design Valutate (Memoria Conversazionale)}

Per la gestione della memoria conversazionale sono state analizzate tre alternative:

\begin{enumerate}
    \item \textbf{Alternativa A -- localStorage} (implementata): history lato client, max 10 messaggi, nessun server-side storage. Semplice e immediata, ma la cronologia è legata al browser.
    \item \textbf{Alternativa B -- Redis}: history lato server con TTL 24h, flush su MongoDB al logout. Più robusta ma complessa da gestire.
    \item \textbf{Alternativa C -- MongoDB}: history permanente sul database, supporto multi-conversazione. Massima persistenza ma overhead significativo.
\end{enumerate}

La scelta dell'Alternativa A è motivata dalla semplicità di implementazione e dall'adeguatezza per il caso d'uso corrente, in linea con il principio AMDD di ``just enough modeling''.

\section{Limitazioni}

\begin{itemize}
    \item \texttt{secure=False} sul cookie (deploy locale senza HTTPS)
    \item Conversazione solo lato client (cambio browser/dispositivo perde la cronologia)
    \item Nessun refresh token (scadenza fissa a 7 giorni)
    \item Nessun accesso a dati reali dell'ateneo
\end{itemize}


% ═══════════════════════════════════════════════════════════════════
% CHAPTER 6 – ITERAZIONE 4
% ═══════════════════════════════════════════════════════════════════
\chapter{Iterazione 4 -- Web Search, Date Esami e Pipeline Logging}

\section{Obiettivo dell'Iterazione}

L'Iterazione 4 arricchisce il sistema con la capacità di recuperare informazioni reali e aggiornate dal sito dell'ateneo. Le macro-funzionalità introdotte sono:
\begin{itemize}
    \item \textbf{QueryAgent}: generazione di query di ricerca web ottimizzate
    \item \textbf{WebAgent}: ricerca web tramite Tavily API su domini \texttt{unibg.it}
    \item \textbf{Categoria \texttt{date\_esami} con routing condizionale}: flusso dedicato per le date degli esami con selezione intelligente del calendario
    \item \textbf{PipelineLogger}: logging strutturato di tutta la pipeline
    \item \textbf{Cambio LLM}: da Groq / LLaMA 3.3 70B a Google Gemini 2.5 Flash
\end{itemize}

\section{User Stories Incluse}

\begin{table}[H]
\centering
\begin{tabularx}{\textwidth}{l X}
\toprule
\textbf{ID} & \textbf{Descrizione} \\
\midrule
US-10 & Recupero informazioni reali dal sito UniBG \\
US-11 & Generazione query di ricerca ottimizzate tenendo conto del profilo \\
US-12 & Tracciamento pipeline in log strutturati \\
US-13 & Risposte con citazione di fonti reali dell'università \\
US-14 & Date esami estratte dal calendario ufficiale del polo \\
\bottomrule
\end{tabularx}
\caption{User stories -- Iterazione 4}
\end{table}

\section{Evoluzione dell'Architettura}

\begin{table}[H]
\centering
\small
\begin{tabularx}{\textwidth}{l X X}
\toprule
\textbf{Funzionalità} & \textbf{Iterazione 3} & \textbf{Iterazione 4} \\
\midrule
Recupero dati reali & Assente (solo LLM) & WebAgent su Tavily \\
Ottimizzazione query & Assente & QueryAgent con Gemini \\
Gestione PDF & Assente & WebAgent + PyMuPDF \\
Logging pipeline & Assente & PipelineLogger su file system \\
LLM Provider & Groq / LLaMA 3.3 70B & Google Gemini 2.5 Flash \\
Workflow & classify → generate → revise & classify → query → [web\_search $|$ exam\_extract] → generate → revise \\
Classificazione & 6 categorie & 7 categorie (+ date\_esami) \\
Routing & Lineare & Condizionale \\
\bottomrule
\end{tabularx}
\caption{Evoluzione dall'Iterazione 3 alla 4}
\end{table}

\section{Workflow LangGraph con Routing Condizionale}

Il workflow LangGraph include ora un \textbf{branch condizionale} dopo il nodo \texttt{query}:

\begin{verbatim}
classify → query → [routing condizionale] → generate → revise → END
                     ├── web_search      (tutte le categorie)
                     └── exam_extract    (solo date_esami)
\end{verbatim}

Il metodo \texttt{\_route\_after\_query(state)} verifica se \texttt{state["category"] == "date\_esami"} e instrada il flusso verso il nodo appropriato. Entrambi i rami convergono su \texttt{generate}.

\section{Query Agent}

Il QueryAgent trasforma la domanda dell'utente in una query di ricerca web breve e ottimizzata (max 6-8 parole chiave), adatta per \texttt{unibg.it}.

\textbf{Regola critica}: il contesto della conversazione ha sempre la priorità sul profilo dello studente. Se la conversazione indica che l'utente sta chiedendo di un corso diverso dal proprio, la query deve riferirsi al corso menzionato nella conversazione.

Input: query + conversation\_history (ultimi 3 turni) + user\_context. Se Gemini fallisce, viene usata la query originale come fallback.

\section{Web Agent}

\subsection{Ricerca Web Standard}
Esegue ricerche su Tavily API limitando i risultati ai domini \texttt{unibg.it} e \texttt{unibg.coursecatalogue.cineca.it}:
\begin{itemize}
    \item Fino a 30 risultati con \texttt{search\_depth="advanced"}
    \item Ordinamento per score decrescente, selezione dei top 5
    \item Estrazione link da PDF tramite PyMuPDF (se disponibile)
    \item Contesto formattato iniettato nel GeneratorAgent
\end{itemize}

\subsection{Gestione Date Esami (\texttt{exam\_extract})}

Quando la categoria è \texttt{date\_esami}, il flusso è:

\begin{enumerate}
    \item \textbf{Selezione calendario}: il dizionario \texttt{EXAM\_CALENDAR\_LINKS} contiene i link ai calendari ufficiali per polo (Ingegneria, Economico-Giuridico, Linguistico, Umanistico) e sessione (invernale, intermedia, estiva, autunnale). L'LLM riceve la domanda e la lista completa e seleziona il calendario appropriato.
    \item \textbf{Regola per ospiti}: se l'utente è ospite e non ha specificato il corso nella domanda, l'LLM restituisce \texttt{NUMERO: 0} indicando l'impossibilità di determinare il calendario.
    \item \textbf{Tavily Extract}: estrazione del contenuto completo dalla URL del calendario scelto.
    \item \textbf{Ricerca aggiuntiva}: Tavily Search per le top 3 fonti web generiche.
    \item \textbf{Formattazione}: contesto combinato (calendario + fonti) passato al GeneratorAgent con prompt specializzato per \texttt{date\_esami}.
\end{enumerate}

\section{Pipeline Logger}

Scrive un file \texttt{.txt} nella cartella \texttt{logs/} per ogni query, con naming convention \texttt{log\_\{matricola\}\_\{timestamp\}.txt}. Ogni file contiene:
\begin{enumerate}
    \item Header (timestamp, status)
    \item Informazioni utente
    \item Query originale e conversation history
    \item Step 1--5 (Classifier, QueryAgent, WebAgent/ExamExtract, Generator, Revision) con system prompt, user prompt e risposta
    \item Tempi di esecuzione per ogni step e totale pipeline
\end{enumerate}

Il logger non blocca mai la pipeline: qualsiasi errore di scrittura viene ingoiato silenziosamente.

\section{Cambio LLM Provider}

\begin{table}[H]
\centering
\begin{tabularx}{\textwidth}{l X X}
\toprule
\textbf{Aspetto} & \textbf{Groq / LLaMA 3.3 70B} & \textbf{Google Gemini 2.5 Flash} \\
\midrule
Context window & 128K token & 1M token \\
Qualità ragionamento & Alta & Molto alta \\
Gestione contesti lunghi & Limitata & Ottima \\
\bottomrule
\end{tabularx}
\caption{Confronto LLM provider}
\end{table}

Il context window più ampio è fondamentale per questa iterazione, dato che il \texttt{web\_context} può contenere testi molto lunghi provenienti da pagine web e PDF.

\section{Classificazione -- Categoria \texttt{date\_esami}}

Il ClassifierAgent è stato esteso con la nuova categoria \texttt{date\_esami}, che intercetta domande relative a date, sessioni e orari degli esami. Le categorie totali sono ora 7: \texttt{informazioni\_corso}, \texttt{orari}, \texttt{date\_esami}, \texttt{procedure}, \texttt{servizi}, \texttt{generale}, \texttt{altro}.

\section{Limitazioni}

\begin{itemize}
    \item Ricerca limitata ai domini \texttt{unibg.it}
    \item Top 5 risultati fissi (nessun re-ranking semantico)
    \item Calendari esami hardcoded (aggiornamento manuale)
    \item Log solo su file system (non consultabili via API)
    \item Nessun refresh token JWT
    \item Conversazione solo in localStorage
    \item Gestione profilo utente non ancora disponibile dal frontend
\end{itemize}


% ═══════════════════════════════════════════════════════════════════
% CHAPTER 7 – ITERAZIONE 5 (WIP)
% ═══════════════════════════════════════════════════════════════════
\chapter{Iterazione 5 -- Gestione Profilo Utente e Miglioramenti (WIP)}

\section{Obiettivo dell'Iterazione}

L'Iterazione 5, attualmente \textit{work in progress}, si concentra su due aree principali:
\begin{itemize}
    \item \textbf{Gestione profilo utente dal frontend}: lo studente autenticato può modificare i propri dati personali (nome, cognome, dipartimento, corso, tipologia, anno) e cambiare la password direttamente dall'interfaccia web
    \item \textbf{Miglioramento del sistema di classificazione}: ottimizzazione del ClassifierAgent per una categorizzazione più precisa delle richieste
\end{itemize}

\section{Funzionalità Implementate (Parziale)}

\subsection{Modifica Profilo Utente}

\textbf{Backend}:
\begin{itemize}
    \item Nuovo endpoint \texttt{PUT /api/auth/profile}: accetta un body con i campi da aggiornare (tutti opzionali). Richiede il cookie JWT per l'autenticazione.
    \item Nuovo endpoint \texttt{PUT /api/auth/password}: accetta \texttt{current\_password} e \texttt{new\_password}. Verifica la password corrente prima di aggiornarla.
    \item Nuovi modelli Pydantic: \texttt{UpdateProfileRequest} (campi opzionali) e \texttt{ChangePasswordRequest}
    \item \texttt{ProfileRepository.updateProfile()}: aggiorna i campi del profilo escludendo \texttt{passwordHash}, \texttt{matricola} e \texttt{\_id} per sicurezza
    \item \texttt{ProfileRepository.updatePassword()}: aggiorna l'hash della password
    \item \texttt{AuthService.updateProfile()} e \texttt{AuthService.changePassword()}: logica applicativa con verifica credenziali
    \item \texttt{AuthController.update\_profile()} e \texttt{AuthController.change\_password()}: orchestrazione e risposta API
\end{itemize}

\textbf{Frontend}:
\begin{itemize}
    \item Nuovo componente \texttt{SettingsPage.jsx}: pagina impostazioni accessibile tramite icona ingranaggio nell'header (visibile solo per utenti autenticati)
    \item Form di modifica profilo con selezione di dipartimento e corso da liste predefinite (stesse mappature di \texttt{LoginPage})
    \item Form di cambio password con verifica della password corrente
    \item Aggiornamento dello stato React (\texttt{userInfo}) dopo la modifica riuscita
\end{itemize}

\subsection{Miglioramento ClassifierAgent}

La categoria \texttt{date\_esami} introdotta nell'Iterazione 4 rappresenta il primo passo verso una classificazione più granulare. L'Iterazione 5 prevede ulteriori ottimizzazioni:
\begin{itemize}
    \item Affinamento dei prompt di classificazione per ridurre le ambiguità
    \item Migliore gestione delle domande follow-up che cambiano contesto
    \item Eventuale introduzione di sotto-categorie per la categoria \texttt{generale}
\end{itemize}

\section{Endpoint API -- Nuovi}

\begin{table}[H]
\centering
\begin{tabularx}{\textwidth}{l l X}
\toprule
\textbf{Metodo} & \textbf{Path} & \textbf{Descrizione} \\
\midrule
PUT & \texttt{/api/auth/profile} & Aggiorna i dati del profilo utente \\
PUT & \texttt{/api/auth/password} & Cambia la password dell'utente \\
\bottomrule
\end{tabularx}
\caption{Nuovi endpoint API -- Iterazione 5}
\end{table}

\section{Stato Attuale}

Questa iterazione è ancora in fase di sviluppo. Le funzionalità di modifica profilo e cambio password sono implementate sia nel backend che nel frontend, ma sono in corso test e raffinamenti dell'interfaccia utente.


% ═══════════════════════════════════════════════════════════════════
% CHAPTER 8 – SCELTE ARCHITETTURALI E DI PROGETTO
% ═══════════════════════════════════════════════════════════════════
\chapter{Scelte Architetturali e di Progetto}

\section{Architettura Multi-Agent con LangGraph}

La scelta fondamentale del progetto è l'adozione di un'architettura \textbf{multi-agent} orchestrata tramite \textbf{LangGraph}. Ogni agente ha una responsabilità specifica e ben definita:

\begin{table}[H]
\centering
\begin{tabularx}{\textwidth}{l X}
\toprule
\textbf{Agente} & \textbf{Responsabilità} \\
\midrule
ClassifierAgent & Classificazione dell'intento in 7 categorie \\
QueryAgent & Generazione di query di ricerca web ottimizzate \\
WebAgent & Ricerca web su Tavily + estrazione calendari esami \\
GeneratorAgent & Generazione della risposta con prompt per categoria \\
RevisionAgent & Revisione per correttezza, struttura, concisione \\
OrchestratorAgent & Coordinamento del workflow completo \\
\bottomrule
\end{tabularx}
\caption{Responsabilità degli agenti}
\end{table}

\textbf{Motivazione}: la separazione in agenti specializzati permette di:
\begin{itemize}
    \item Modificare o sostituire un singolo agente senza impattare gli altri
    \item Aggiungere nuovi nodi alla pipeline (es.\ QueryAgent, WebAgent) senza riscrivere la logica esistente
    \item Testare e debuggare ogni componente in isolamento
    \item Implementare routing condizionale (es.\ exam\_extract) per flussi specializzati
\end{itemize}

\section{Pattern AgentState (Shared State)}

Lo stato della pipeline è gestito tramite un \texttt{TypedDict} (\texttt{AgentState}) condiviso tra tutti i nodi del grafo LangGraph. Ogni nodo legge e scrive i campi di cui ha bisogno, garantendo:
\begin{itemize}
    \item Tracciabilità completa del percorso (campo \texttt{workflow\_steps})
    \item Fallback robusti (se un agente fallisce, il campo non viene sovrascritto e l'agente successivo usa i dati disponibili)
    \item Estensibilità: nuovi campi possono essere aggiunti senza modificare i nodi esistenti
\end{itemize}

\section{Scelta del LLM Provider}

Il progetto ha attraversato due scelte di LLM provider:
\begin{enumerate}
    \item \textbf{Iterazioni 1--3}: Groq con LLaMA 3.3 70B Versatile. Scelto per la latenza molto bassa e la buona qualità delle risposte.
    \item \textbf{Iterazione 4+}: Google Gemini 2.5 Flash. Motivato dal context window di 1M token (vs 128K), necessario per iniettare i contenuti web nel prompt, e dalla qualità superiore nel ragionamento.
\end{enumerate}

L'architettura è stata progettata per rendere il cambio di provider \textbf{trasparente}: l'LLM è iniettato dall'OrchestratorAgent come dependency e nessun agente contiene riferimenti hardcoded al modello.

\section{Autenticazione con httpOnly Cookie}

La scelta di trasmettere il JWT tramite \textbf{httpOnly cookie} (anziché in un header \texttt{Authorization}) è motivata da:
\begin{itemize}
    \item \textbf{Protezione XSS}: il token non è accessibile da JavaScript, riducendo il rischio di furto
    \item \textbf{Semplicità frontend}: il browser allega automaticamente il cookie ad ogni richiesta con \texttt{credentials: "include"}
    \item \textbf{SameSite=lax}: protezione CSRF di base senza necessità di token CSRF espliciti
\end{itemize}

\section{Memoria Conversazionale lato Client}

La scelta di mantenere la conversation history in \texttt{localStorage} (anziché su Redis o MongoDB) è motivata da:
\begin{itemize}
    \item Semplicità di implementazione
    \item Nessun overhead di infrastruttura (Redis) o latenza (MongoDB)
    \item Sufficiente per il caso d'uso corrente (singolo browser/dispositivo)
    \item Limitazione a 10 messaggi per contenere il context window
\end{itemize}

\section{Tavily API per la Ricerca Web}

Tavily è stato scelto come motore di ricerca web per:
\begin{itemize}
    \item API semplice con filtering per dominio (\texttt{include\_domains})
    \item Supporto \texttt{search\_depth="advanced"} per contenuti più ricchi
    \item Funzionalità \texttt{extract()} per estrazione contenuto da URL specifiche
    \item Restituzione di \texttt{raw\_content} per avere il testo completo delle pagine
\end{itemize}

\section{Routing Condizionale per Date Esami}

La scelta di un flusso dedicato per la categoria \texttt{date\_esami} è motivata dalla specificità del dominio:
\begin{itemize}
    \item I calendari esami sono organizzati per polo/dipartimento e sessione con URL specifiche
    \item Una ricerca web generica non è sufficiente per recuperare date precise di un esame specifico
    \item Il flusso \texttt{exam\_extract} usa Tavily Extract (non search) per estrarre il contenuto completo del calendario
    \item L'LLM seleziona autonomamente il calendario corretto in base alla domanda e al profilo
\end{itemize}

\section{Scelte di Sicurezza}

\begin{itemize}
    \item \textbf{Password}: hashate con bcrypt (salt incluso), mai salvate in chiaro
    \item \textbf{JWT}: firma HMAC-SHA256, scadenza 7 giorni, trasmissione via httpOnly cookie
    \item \textbf{CORS}: origini esplicite in produzione, \texttt{allow\_credentials=True}
    \item \textbf{Public Profile}: il campo \texttt{passwordHash} è filtrato da ogni risposta API
    \item \textbf{Update Profile}: i campi \texttt{passwordHash}, \texttt{matricola} e \texttt{\_id} sono esclusi dagli aggiornamenti per sicurezza
\end{itemize}

\section{Containerizzazione}

Docker Compose permette il deploy dell'intero stack con un singolo comando (\texttt{docker-compose up}):
\begin{itemize}
    \item Il backend è isolato in un container Python
    \item Il frontend è servito da Nginx con build Vite
    \item MongoDB Atlas è esterno ai container (SaaS)
    \item I volumi permettono hot-reload in sviluppo
\end{itemize}


% ═══════════════════════════════════════════════════════════════════
% CHAPTER 9 – CONCLUSIONI
% ═══════════════════════════════════════════════════════════════════
\chapter{Conclusioni e Sviluppi Futuri}

\section{Riepilogo delle Iterazioni}

\begin{table}[H]
\centering
\small
\begin{tabularx}{\textwidth}{l X}
\toprule
\textbf{Iterazione} & \textbf{Funzionalità Principali} \\
\midrule
0 -- Visione & Requisiti, user stories, architettura target \\
1 -- MVP & Pipeline classify → generate → revise (Groq/LLaMA) \\
2 -- Auth & Registrazione, login, MongoDB, personalizzazione, React, Docker \\
3 -- Sessione & JWT httpOnly cookie, verify/logout, localStorage memory, auth refactoring \\
4 -- Web \& Esami & QueryAgent, WebAgent (Tavily), date\_esami con routing condizionale, PipelineLogger, Gemini 2.5 Flash \\
5 -- Profilo (WIP) & Modifica profilo utente, cambio password, miglioramenti classifier \\
\bottomrule
\end{tabularx}
\caption{Riepilogo iterazioni}
\end{table}

\section{Obiettivi Raggiunti}

Il sistema Agentic UniBG ha raggiunto i seguenti obiettivi:
\begin{itemize}
    \item Pipeline multi-agent funzionante con 5 agenti specializzati
    \item Classificazione automatica degli intenti in 7 categorie
    \item Recupero informazioni reali dal sito \texttt{unibg.it} tramite Tavily
    \item Flusso dedicato per date esami con selezione intelligente del calendario
    \item Autenticazione sicura con JWT httpOnly cookie
    \item Personalizzazione delle risposte basata sul profilo studente
    \item Memoria conversazionale per risposte coerenti con il contesto
    \item Logging completo della pipeline per diagnostica e audit
    \item Deploy containerizzato con Docker Compose
\end{itemize}

\section{Sviluppi Futuri}

\begin{itemize}
    \item \textbf{Completamento Iterazione 5}: finalizzazione della gestione profilo utente e ottimizzazione del ClassifierAgent
    \item \textbf{Refresh token}: implementazione di un meccanismo di rinnovo del JWT per evitare ri-autenticazioni dopo 7 giorni
    \item \textbf{Persistenza conversazione su backend}: migrazione da localStorage a MongoDB per supporto multi-dispositivo
    \item \textbf{Re-ranking semantico}: miglioramento della selezione dei risultati web mediante similarità semantica
    \item \textbf{Log consultabili via API}: esposizione dei log della pipeline tramite endpoint REST per facilitare il monitoraggio
    \item \textbf{Deploy in produzione}: configurazione HTTPS, \texttt{secure=True} sul cookie, dominio personalizzato
    \item \textbf{Aggiornamento automatico calendari}: meccanismo per mantenere aggiornati i link ai calendari esami senza intervento manuale
\end{itemize}

\section{Considerazioni sulla Metodologia AMDD}

L'approccio AMDD si è rivelato efficace per questo progetto:
\begin{itemize}
    \item La visione iniziale (Iterazione 0) ha fornito una direzione chiara senza vincolare le scelte implementative
    \item Ogni iterazione ha prodotto un incremento funzionante e testabile
    \item I modelli UML sono stati creati e aggiornati ``just enough'' per documentare le scelte senza overhead
    \item Le decisioni architetturali (es.\ cambio da Groq a Gemini, routing condizionale per esami) sono emerse naturalmente durante le iterazioni
    \item L'approccio incrementale ha permesso di scoprire e risolvere problemi (es.\ necessità del JWT, limiti del context window) senza rework massivo
\end{itemize}

\end{document}
